\subsection{\textbf{Broadband Sensor Noise and Channel Dropout}}
\newcommand{\CurrentNoiseTable}{%
\begin{table*}[!t]
\caption{Frozen-artifact accuracy under broadband sensor noise and channel dropout.}
\label{tab:noise}
\centering\footnotesize
\setlength{\tabcolsep}{3pt}
\begin{tabular}{lrrrrr}
\toprule
Noise & Official RMSE & Official MAE & Cov. (\%) & Internal RMSE & Cov. (\%) \\
\midrule
none & 12.01 & 8.52 & 99.3 & 13.30 & 99.0 \\
5\% & 12.05 & 8.60 & 99.3 & 13.37 & 99.0 \\
10\% & 12.17 & 8.62 & 99.2 & 13.55 & 99.2 \\
20\% & 12.37 & 8.76 & 98.9 & 13.84 & 98.8 \\
30\% & 12.67 & 9.04 & 99.0 & 14.29 & 98.8 \\
50\% & 13.68 & 9.97 & 99.2 & 15.44 & 97.9 \\
2\% cells & 11.96 & 8.56 & 99.6 & 13.37 & 99.0 \\
5\% cells & 12.23 & 8.77 & 99.2 & 13.54 & 99.0 \\
\bottomrule
\end{tabular}
\par\vspace{2pt}
\begin{minipage}{\textwidth}\scriptsize Rows 1--6 add per-channel Gaussian noise; the level is a multiple of the fitting-only within-regime sensor scale. Rows 7--8 instead drop whole cells at zero added noise. Settings and speed commands are unperturbed. Noise is applied after model freeze; no weight, threshold, corrector or interval radius was fitted on these levels. Coverage is empirical clean endpoint/prefix coverage under the frozen radius, not a new guarantee.\end{minipage}
\end{table*}
}

\CurrentNoiseTable
The structured protocol models failures confined to related channels. Real
acquisition also carries broadband noise on every channel at once and occasional
missing cells. We therefore add an independent measurement layer that perturbs
raw sensor readings before the frozen fitting-only scaler: per-channel Gaussian
noise scaled to the within-regime sensor scale, plus independent cell dropouts.
Settings and speed commands stay untouched, and the layer is applied only after
freeze, so no weight, detector threshold, corrector or interval radius was fitted
on it. Levels are nested: one per-engine draw is scaled, so a sweep is paired
rather than resampled.

Table~\ref{tab:noise} shows graceful degradation. At 20\% of the within-regime
scale, pooled official RMSE moves from \CurrentNoiseOfficialZero{} to
\CurrentNoiseOfficialTwenty{} cycles, an increase of 3.0\%, and empirical coverage
remains \CurrentNoiseCoverageTwenty{}\%. At 50\% it is
\CurrentNoiseOfficialFifty{} cycles with \CurrentNoiseCoverageFifty{}\% coverage.
The internal partition follows the same pattern, from \CurrentNoiseHeldoutZero{}
to \CurrentNoiseHeldoutFifty{}. Modest dropout is not harmful: dropping 2\% of
cells leaves pooled official RMSE at \CurrentNoiseDropoutTwo{}, and 5\% gives
\CurrentNoiseDropoutFive{}. These are sensitivity measurements of one frozen
artifact, not a claim that any level matches a measured sensor specification.

\subsection{\textbf{Comparison with the Supplied Studies}}
\newcommand{\CurrentLiteratureTable}{%
\begin{table*}[!t]
\caption{Reported C-MAPSS results at the 125-cycle target against the supplied studies.}
\label{tab:literature}
\centering\footnotesize
\setlength{\tabcolsep}{3pt}
\begin{tabular}{lrrrr}
\toprule
Model & FD001 & FD002 & FD003 & FD004 \\
\midrule
GRACE-RUL (one shared model, no subset label) & \textbf{11.93} & \textbf{11.77} & \textbf{11.25} & \textbf{12.57} \\
Hybrid CNN--GRU (reported, five repeats) & 13.1 & 16.8 & 14.5 & 18.9 \\
TRANS-D3 (reported) & 1.66 & 13.23 & 6.61 & 12.26 \\
Best third-party baseline (reported) & 10.61 (STAR) & 13.47 (STAR) & 10.71 (STAR) & 15.86 (DLformer) \\
Reported value ahead of GRACE-RUL & TRANS-D3 1.66 & none & TRANS-D3 6.61 & TRANS-D3 12.26 \\
Position among third-party baselines (proposed row excluded) & 6 & 1 & 3 & 1 \\
\bottomrule
\end{tabular}
\par\vspace{2pt}
\begin{minipage}{\textwidth}\scriptsize Reported columns are transcribed from the cited papers and are not recomputed here. Rank counts only third-party baselines, so a proposed row cannot define its own benchmark. The penultimate row names every reported value still ahead of this work, including the TRANS-D3 proposed row where it applies; on FD004 this work is second overall, not first. This work is one shared model without a subset label, whereas the compared entries are subset-specific. Official C-MAPSS tests were inspected in earlier development, so all values are exploratory rather than a pristine holdout.\end{minipage}
\end{table*}
}

\CurrentLiteratureTable
Table~\ref{tab:literature} transcribes the reported values of the supplied
studies. This work is first on FD002 among every reported row, second on FD004
behind the reported TRANS-D3 value of 12.26, third on FD003 among third-party
baselines, and \CurrentRankFdA{}th on FD001. It is not first on FD004 and no
such claim is made. The hybrid CNN--GRU study reports
13.1, 16.8, 14.5 and 18.9, so this model improves every subset by 1.17 to 6.33
cycles at a comparable 125-cycle target.

Two qualifications matter more than the ordering. First, the compared entries are
subset-specific models, whereas this work uses one shared model without a subset
label, so the comparison is not budget-matched in the harder direction. Second,
the TRANS-D3 FD001 and FD003 values of 1.66 and 6.61 lie far below every
third-party entry in the same table, whose best values are \CurrentBestFdA{}
and \CurrentBestFdB{}; no capped C-MAPSS protocol in this study reaches them, so they are
retained for completeness rather than treated as a reproducible target. The
partial-failure study reports a single pooled fault-conditioned RMSE of
\CurrentPaperBOfficial{} with unspecified channel coverage and severity, which is
not decomposable by subset and is therefore not entered as a per-subset cell.
